% !PRINTED 532
\weberpage{532}
But the following theorem also holds with respect to the \gesperrt{functionals}:

2. \gesperrt{If $\omega_1, \omega_2, \ldots, \omega_n$ is a basis of $\mathfrak{o}$, then in the form}
\weq{9}{\omega = u_1\omega_1 + u_2\omega_2 + \cdots + u_n\omega_n,}
\gesperrt{in which $u_1, u_2, \ldots, u_n$ are integral rational functionals, all integral functionals in $\Omega$ are contained.}

For we have already shown above that all functionals whatsoever are contained in the form (9) when the coefficients $u_1, u_2, \ldots, u_n$ are integral or fractional rational functionals. But we can always determine an integral primitive function $e$ in such a way that
\[
e u_1 = y_1, \; e u_2 = y_2, \; \ldots, \; e u_n = y_n
\]
are integral functions of the variables, and then
\weq{10}{e\omega = y_1\omega_1 + y_2\omega_2 + \cdots + y_n\omega_n,}
and since $e$ is a unit, $e\omega$ is integral together with $\omega$. Now since the coefficients of the powers and products of the variables in the function (10) must, by \S.\,142, 3., be integers, it follows by 1. that the coefficients in the functions $y_1, y_2, \ldots, y_n$ must be integral rational numbers, and that consequently $u_1, u_2, \ldots, u_n$ are integral rational functionals.

\begin{center}
\S.\,146.

\gesperrt{The Bases of the Functionals.}
\end{center}

\gesperrt{If $\mu$ is an integral functional of the field $\Omega$, then by a basis of $\mu$ we understand a system of $n$ integers in $\Omega$:}
\weq{1}{\alpha_1, \alpha_2, \ldots, \alpha_n,}
\gesperrt{which is a basis of the field $\Omega$, and which has the property that in the form}
\weq{2}{\alpha = x_1\alpha_1 + x_2\alpha_2 + \cdots + x_n\alpha_n}
\gesperrt{all the integers of the field $\Omega$ that are divisible by $\mu$, and no others, are contained, when integral rational numbers are put for $x_1, x_2, \ldots, x_n$.}

It shall now be proved that every integral functional has a basis.
